% ---------------------------------------------------------------------
%  CHƯƠNG 5 — KẾT QUẢ VÀ ĐÁNH GIÁ
% ---------------------------------------------------------------------
\chapter{Kết quả và đánh giá}
\label{ch:ketqua}

\begin{muctieu}
Chương này trình bày kết quả đạt được, các độ đo đánh giá định lượng, phần
bàn luận và những hạn chế còn tồn tại của hệ thống gợi ý phòng trọ.
\end{muctieu}

\section{Kết quả đạt được}

Hệ thống đã hoàn thành các mục tiêu đề ra ở Chương~\ref{ch:modau}. Bộ dữ
liệu 2.500 phòng trọ tại sáu quận trung tâm Đà Nẵng đã được xây dựng và làm
sạch, mô hình gợi ý dựa trên nội dung đạt Precision@5 bằng 0,76, vượt
ngưỡng mục tiêu 0,7, và giao diện web cho phép người dùng nhập tiêu chí và
nhận danh sách phòng trọ gợi ý kèm bản đồ trong thời gian thực.

\section{Độ đo đánh giá}

Vì bài toán là xếp hạng và gợi ý chứ không phải phân lớp, đồ án dùng các độ
đo phổ biến cho hệ gợi ý gồm Precision@5, Recall@5 và NDCG@5, tính trên
lịch sử tương tác đã lưu hoặc đã liên hệ của 200 người dùng thử nghiệm, đối
chiếu với 500 phòng trọ ở tập kiểm thử ngoại tuyến. Bảng~\ref{tab:ketqua}
so sánh mô hình chỉ dùng lọc dựa trên nội dung với phiên bản kết hợp thêm
thành phần cộng tác đơn giản.

\begin{table}[H]
\centering
\caption{So sánh hiệu năng các phiên bản mô hình gợi ý.}
\label{tab:ketqua}
\begin{tabular*}{\textwidth}{@{\extracolsep{\fill}}>{\bfseries\color{daunavy}}l ccc}
\toprule
\rowcolor{gray100}\bfseries\color{daunavy}Mô hình & \bfseries\color{daumaroon}Precision@5 & \bfseries\color{daumaroon}Recall@5 & \bfseries\color{daumaroon}NDCG@5 \\
\midrule
Chỉ lọc dựa trên nội dung          & 0,69 & 0,60 & 0,75 \\
Kết hợp thành phần cộng tác        & 0,76 & 0,68 & 0,81 \\
\bottomrule
\end{tabular*}
\end{table}

Ngoài đánh giá ngoại tuyến, đồ án tổ chức khảo sát nhanh trên 30 sinh viên
dùng thử hệ thống, chấm mức độ hài lòng theo thang điểm 5. Kết quả trung
bình đạt 4,1 điểm, trong đó tiêu chí được đánh giá cao nhất là tốc độ phản
hồi và tiêu chí thấp nhất là độ đa dạng của kết quả gợi ý.

\section{Bàn luận}

Phiên bản kết hợp thành phần cộng tác nhỉnh hơn phiên bản chỉ dùng nội dung
khoảng bảy điểm phần trăm ở Precision@5. Mức cải thiện này đến từ việc một
số phòng trọ tốt nhưng có mô tả sơ sài, dẫn đến độ tương đồng văn bản thấp,
vẫn được gợi ý nhờ được người dùng có nhu cầu tương tự lưu lại trước đó.
Phân tích các trường hợp gợi ý chưa tốt cho thấy phần lớn tập trung ở những
tiêu chí có khu vực mong muốn ít phòng trọ, khi đó hệ thống buộc phải nới
lỏng điều kiện khu vực để trả đủ năm kết quả, làm giảm độ phù hợp trung
bình của danh sách gợi ý.

\section{Hạn chế}

\begin{luuy}
Hệ thống còn gợi ý kém đa dạng khi tiêu chí người dùng quá chặt, do bản chất
lọc dựa trên nội dung có xu hướng trả về những phòng trọ giống nhau. Thành
phần cộng tác chỉ dựa trên năm người dùng gần nhất nên chưa phát huy hết
tác dụng khi số lượng người dùng còn ít. Dữ liệu phòng trọ cũng chưa được
cập nhật tự động, một số tin đăng có thể đã hết phòng nhưng vẫn còn trong
hệ thống.
\end{luuy}
